\SourceLeaf{266}
\fancyhead[L]{\small UNTERSUCHUNGEN ÜBER THEMEN DER HÖHEREN GEODÄSIE}
\renewcommand{\Article}[1]{\par\medskip\Needspace{4\baselineskip}\pdfbookmark[1]{Art. #1}{geodaesie1-art-#1}\centerline{\textbf{#1.}}\nobreak\smallskip}
\begin{center}
UNTERSUCHUNGEN\par\bigskip
ÜBER\par\bigskip
{\large THEMEN DER HÖHEREN GEODÄSIE.}\par\medskip
\rule{.65in}{.3pt}
\end{center}
\noindent\hspace*{\parindent}Bei den trigonometrischen Vermessungen, die sich über das ganze Königreich Hannover erstreckten und die ich teils selbst durchführte, teils leitete, wurden Wege eingeschlagen, die von den sonst üblichen abweichen. Das betrifft die Art der Messungen und noch mehr ihre anschließende mathematische Behandlung und die Gewinnung der Ergebnisse. Ursprünglich hatte ich vor, nach dem vollständigen Abschluss der Messungen diese samt allen von mir angewandten Verfahren in einem eigenen Werk darzustellen. Aus Gründen, die hier nicht zur Sache gehören, konnte ich diesen Plan bisher nicht ausführen. Daher wähle ich den Weg, das Besondere an meiner theoretischen Arbeit in einer Reihe von Abhandlungen bekannt zu machen. Das tue ich umso lieber, als ich so die Freiheit behalte, manche Untersuchungen ausführlich zu entwickeln, die für sich genommen von Interesse sind und mit den übrigen in engem Zusammenhang stehen, auch wenn ich sie bei meinen Vermessungen nicht unmittelbar angewandt habe. Das gilt insbesondere für den größten Teil der vorliegenden ersten Abhandlung.
